\bookchapter{Κεφάλαιο 2: Οργανώστε το ψηφιακό σας δω$\mu$άτιο — Διαχείριση αρχείων}

\chapterintro{Το κινητό$\mu$ε παρατετα$\mu$ένη χρήση ακαταστατεί, η καλή συνήθεια ξεκινά από τη διαχείριση αρχείων. Όπως το δω$\mu$άτιό σας χρειάζεται περιοδική οργάνωση, έτσι και ο ψηφιακός σας χώρος χρειάζεται φροντίδα. Ένα καλά οργανω$\mu$ένο κινητό, όχι $\mu$όνο βρίσκετε εύκολα τα πράγ$\mu$ατα, αλλά λειτουργεί πιο ο$\mu$αλά.}

\booksection{I. Σύστη$\mu$α αποθήκευσης τριών επιπέδων: το ψηφιακό σας αποθετήριο}

\begin{enumerate}
  \item \textbf{Εσωτερική αποθήκευση}: Η βασική ζώνη, γρήγορη, εδώ βρίσκονται το σύστη$\mu$α και οι εφαρ$\mu$ογές. Ο πιο πολύτι$\mu$ος χώρος του τηλεφώνου σας, πρέπει να τον διαχειρίζεστε προσεκτικά.
  \item \textbf{SD κάρτα (αν υποστηρίζεται)}: επέκταση ντουλάπας, λίγο πιο αργή, κατάλληλη για $\mu$εγάλα αρχεία όπως ταινίες, φωτογραφίες. Προσοχή: πολλές ση$\mu$ερινές ναυαρχίδες δεν υποστηρίζουν πλέον κάρτες SD.
  \item \textbf{Cloud αποθήκευση}: Χρη$\mu$ατοκιβώτιο + συγχρονισ$\mu$ένος cloud δίσκος. Προστασία από απώλεια και συγχρονισ$\mu$ός $\mu$εταξύ πολλαπλών συσκευών. iCloud, Baidu Netdisk, OneDrive είναι κοινές επιλογές.
\end{enumerate}

\booksection{II. Γνωρίστε τον «διαχειριστή αρχείων»}

Το Android συνήθως έχει $\mu$ια [Αρχεία]ή [Διαχείριση Αρχείων] App, το iOS έχει την ενσω$\mu$ατω$\mu$ένη [Αρχεία] App. Είναι ο «εξευρευτής πόρων» για να βλέπετε όλα τα αρχεία του κινητού.

\booksubsection{Ση$\mu$ασία κοινών φακέλων:}

\begin{itemize}
  \item \textbf{DCIM}: φωτογραφίες και βίντεο από την κά$\mu$ερα.
  \item \textbf{Pictures}: άλλες εικόνες, όπως στιγ$\mu$ιότυπα οθόνης, αποθηκευ$\mu$ένες εικόνες.
  \item \textbf{Download}: αρχεία λήψεων browser.
  \item \textbf{Music}: αρχείο $\mu$ουσικής.
  \item \textbf{Movies}: αρχείο βίντεο.
  \item \textbf{Documents}: αρχεία εγγράφων.
  \item \textbf{Android}: φάκελος συστή$\mu$ατος Android, περιέχει δεδο$\mu$ένα κάθε εφαρ$\mu$ογής, γενικά $\mu$ην τα διαγράφετε χειροκίνητα.
\end{itemize}

\booksection{III. Δο$\mu$ή συστή$\mu$ατος αρχείων: βαθύτερη κατανόηση}

Το Android βασίζεται σε Linux, το σύστη$\mu$α αρχείων έχει$\mu$η-επίπεδη δο$\mu$ή. Αν και οικανονικό χρήστες δεν χρειάζεται να κατανοήσουν σε βάθος, η γνώση$\mu$ερικά βασικών έννοιες βοηθάει.

\booksubsection{Δια$\mu$ερισ$\mu$ατοποίηση αποθήκευσης:}

\begin{itemize}
  \item \textbf{/system}: Δια$\mu$έρισ$\mu$α συστή$\mu$ατος, περιέχει τα αρχεία του λειτουργικού συστή$\mu$ατος. Οι απλοί χρήστες δεν έχουν πρόσβαση, απαιτούνται δι και ώ$\mu$ατα root.
  \item \textbf{/data}: δια$\mu$έρισ$\mu$α δεδο$\mu$ένων, αποθηκεύει δεδο$\mu$ένα χρήστη και εφαρ$\mu$ογών. Αυτός είναι ο καθη$\mu$ερινός σας χώρος.
  \item \textbf{/sdcard ή /storage/emulated/0}: ο αποθηκευτικός χώρος προσβάσι$\mu$ος από τον χρήστη, οι φωτογραφίες σας, τα ληφθέντα αρχεία είναι όλα εδώ.
\end{itemize}

\booksubsection{Αποθήκευση δεδο$\mu$ένων App:}

\begin{itemize}
  \item \textbf{Εσωτερική αποθήκευση}: Ο ιδιωτικός χώρος της εφαρ$\mu$ογής, απρόσιτος από άλλες εφαρ$\mu$ογές. Διαγράφεται $\mu$ε την απεγκατάσταση.
  \item \textbf{Εξωτερική αποθήκευση}: αρχεία που αποθηκεύονται από την εφαρ$\mu$ογή σε περιοχή προσβάσι$\mu$η από τον χρήστη, όπως φωτογραφίες, λήψεις. Μπορεί να $\mu$ην διαγραφούν κατά την απεγκατάσταση της εφαρ$\mu$ογής.
\end{itemize}

\booksection{IV. Αντίγραφο ασφαλείας και ανάκτηση δεδο$\mu$ένων: ασφαλίστε τα δεδο$\mu$ένα σας}

Τα δεδο$\mu$ένα δεν έχουν τι$\mu$ή, το αντίγραφο ασφαλείαςείναι απαραίτητο συνήθεια.

\booksubsection{Στρατηγική αντίγραφου ασφαλείας:}

\begin{enumerate}
  \item \textbf{Cloud backup}: Αυτό$\mu$ατα, βολικά, αλλά απαιτούν δίκτυο και χώρο cloud. iCloud, Huawei Cloud, Xiaomi Cloud το υποστηρίζουν.
  \item \textbf{Τοπικό αντίγραφο}: Συνδέστε τον υπολογιστή $\mu$ε καλώδιο, χρησι$\mu$οποιήστε το επίση$\mu$ο εργαλείο backup. Γρήγορο, $\mu$εγάλος χώρος, αλλά απαιτεί χειροκίνητη λειτουργία.
  \item \textbf{Διπλό αντίγραφο}: cloud αντίγραφο ασφαλείας + τοπικό αντίγραφο ασφαλείας, το πιο ασφαλές.
\end{enumerate}

\booksubsection{Τι να αντιγράψετε;}

\begin{itemize}
  \item \textbf{Φωτογραφίες και βίντεο}: καταλα$\mu$βάνει τον περισσότερο χώρο, αλλά είναι και το πολυτι$\mu$ότερο.
  \item \textbf{Επαφές}: Πιο βολικό στον cloud λογαριασ$\mu$ό.
  \item \textbf{Ιστορικό συνο$\mu$ιλιών WeChat}: Χρησι$\mu$οποιήστε την ενσω$\mu$ατω$\mu$ένη λειτουργία backup του WeChat.
  \item \textbf{Ση$\mu$αντικά έγγραφα}: αποθηκεύστε στο cloud ή στον υπολογιστή.
\end{itemize}

\booksection{V. Ασφάλεια αρχείων: προστατεύστε την ιδιωτικότητά σας}

\begin{itemize}
  \item \textbf{Κρυπτογράφηση ευαίσθητων αρχείων}: πολλά κινητά έχουν λειτουργία "ασφαλούς φακέλου " ή "ιδιωτικού χώρου", $\mu$πορείτε να βάλετε ευαίσθητα αρχεία εκεί, χρειάζεται κωδικός ή δακτυλικό αποτύπω$\mu$α για πρόσβαση.
  \item \textbf{Διαγραφή δεν ση$\mu$αίνει καταστροφή}: η απλή διαγραφή απλώς επιση$\mu$αίνει τον χώρο ως επανεγγράψι$\mu$ο, τα δεδο$\mu$ένα παρα$\mu$ένουν. Για πλήρη διαγραφή, χρειάζεται ειδικό εργαλείο ασφαλούς διαγραφής ή πολλαπλή επανεγγραφή.
  \item \textbf{Cloud ασφάλεια}: επιλέξτε αξιόπιστο πάροχο cloud, ενεργοποιήστε τον έλεγχο ταυτότητας δύο παραγόντων, ελέγχετε τακτικά τις συνδεδε$\mu$ένες συσκευές.
\end{itemize}

\booksection{VI. Δεν επαρκεί ο χώρος; Κάντε καθαρισ$\mu$ό}

\begin{itemize}
  \item \textbf{Εκκαθάριση προσωρινής $\mu$νή$\mu$ης}: η προσωρινή $\mu$νή$\mu$η της εφαρ$\mu$ογής, καθαρίζοντάς την ελευθερώνετε χώρο. Μπορείτε να την καθαρίσετε στις Ρυθ$\mu$ίσεις > Αποθήκευση.
  \item \textbf{Εύρεση $\mu$εγάλων αρχείων}: ταξινο$\mu$ήστε αρχεία κατά $\mu$έγεθος, βρείτε τους «ενόχους» που καταλα$\mu$βάνουν χώρο.
  \item \textbf{Διαγραφή υπολει$\mu$$\mu$άτων}: $\mu$ετά την απεγκατάσταση $\mu$ιας εφαρ$\mu$ογής, $\mu$ερικές φορές $\mu$ένουν αρχεία σκουπιδιών, καθαρίστε τα χειροκίνητα.
  \item \textbf{Καθάρισ$\mu$α άλ$\mu$που$\mu$}: διαγραφή διπλότυπων, θολών, περιττών screenshots.
  \item \textbf{Καθάρισ$\mu$α WeChat}: το WeChat είναι $\mu$εγάλος καταναλωτής χώρου, ο τακτικός καθαρισ$\mu$ός ιστορικού συνο$\mu$ιλιών και προσωρινής $\mu$νή$\mu$ης είναι απαραίτητος.
\end{itemize}

\booksection{VII. Πλήρης ανάλυση κοινών $\mu$ορφών αρχείων}

Οι διάφορες καταλήξεις αρχείων στο κινητό $\mu$περδεύουν, αν τις καταλάβετε $\mu$πορεί να διαχειριστείτε καλύτερα τα αρχεία:

\begin{center}
\begin{tabularx}{\linewidth}{|X|X|X|X|}
\hline
\textbf{Μορφή} & \textbf{Τύπος} & \textbf{Χαρακτηριστικά} & \textbf{Προτεινό$\mu$ενος τρόπος ανοίγ$\mu$ατος} \tabularnewline
\hline
JPG/JPEG & Εικόνα & Lossy συ$\mu$πίεση, $\mu$ικρό $\mu$έγεθος, καλύτερη συ$\mu$βατότητα. Προεπιλεγ$\mu$ένη $\mu$ορφή φωτογραφίας κινητού. & Άλ$\mu$που$\mu$/οποιοδήποτε image viewer \tabularnewline
PNG & Εικόνα & Lossless συ$\mu$πίεση, υποστηρίζει διαφανές φόντο. Σε screenshots, εικονίδιασυχνά χρησι$\mu$οποιού$\mu$ενα. Μεγαλύτερο από JPG. & Άλ$\mu$που$\mu$/οποιοδήποτε image viewer \tabularnewline
MP3 & Ήχος & Lossy συ$\mu$πίεση, $\mu$ικρό $\mu$έγεθος, εξαιρετική συ$\mu$βατότητα. Πιο κοινή $\mu$ορφή ήχου. & Όποιο music player \tabularnewline
FLAC & Ήχος & Lossless συ$\mu$πίεση, ποιότητα ίδια $\mu$ε CD, $\mu$έγεθος 3-5 φορές $\mu$εγαλύτερο από MP3. Προτί$\mu$ηση audiophiles. & Music player/επαγγελ$\mu$ατική App \tabularnewline
WAV & Ήχος & Ασυ$\mu$πίεστος αρχικό ήχος, καλύτερη ποιότητα αλλά τεράστιο $\mu$έγεθος. Σε επαγγελ$\mu$ατική ηχογράφηση συχνά χρησι$\mu$οποιού$\mu$ενα. & Music player \tabularnewline
M4A & Ήχος & Κωδικοποίηση AAC, καλύτερη ποιότητα από MP3 στο ίδιο bitrate. Στο οικοσύστη$\mu$α Apple συχνά χρησι$\mu$οποιού$\mu$ενα. & Apple Music/Music player \tabularnewline
MP4 & Βίντεο & Πιο κοινό container format βίντεο, εξαιρετική συ$\mu$βατότητα. Προεπιλογή βίντεο κινητού. & Άλ$\mu$που$\mu$/οποιοδήποτε video player \tabularnewline
MKV & Βίντεο & Ανοιχτού κώδικα container, υποστηρίζει πολλαπλάκο$\mu$$\mu$άια ήχου και υποτίτλους. Σε πόρους ταινιώνσυχνά χρησι$\mu$οποιού$\mu$ενα. & VLC/nPlayer \tabularnewline
DOC/DOCX & Έγγραφα & Έγγραφο Word. Το DOCX είναι συ$\mu$πιεσ$\mu$ένη $\mu$ορφή, $\mu$ικρότερο $\mu$έγεθος. Το DOC είναι παλιά $\mu$ορφή. & WPS/Office \tabularnewline
PPT/PPTX & Παρουσίαση & PowerPoint παρουσίαση. Το PPTX είναι συ$\mu$πιεσ$\mu$ένη $\mu$ορφή. & WPS/Office \tabularnewline
XLS/XLSX & Πίνακας & Excel λογιστικό φύλλο. Το XLSX είναι συ$\mu$πιεσ$\mu$ένη $\mu$ορφή. & WPS/Office \tabularnewline
PDF & Έγγραφα & Cross-platform $\mu$ορφή, σταθερή διάταξη $\mu$ηεπεξεργάσι$\mu$ο. Σε συ$\mu$βόλαια, διατριβές, ebooks συχνά χρησι$\mu$οποιού$\mu$ενα. & Adobe Reader/WPS \tabularnewline
ZIP & Συ$\mu$πίεση & Πιο κοινή $\mu$ορφή συ$\mu$πίεσης, cross-platform συ$\mu$βατότητα. & Ενσω$\mu$ατω$\mu$ένο/αποσυ$\mu$πίεση App \tabularnewline
RAR & Συ$\mu$πίεση & Μεγαλύτερος λόγος συ$\mu$πίεσης από ZIP, αλλά χρειάζεται ειδικό λογισ$\mu$ικό. & αποσυ$\mu$πίεση App \tabularnewline
ISO & Εικόνα & Αρχείο εικόνας δίσκου, πλήρης αντιγραφή περιεχο$\mu$ένου δίσκου. Σε πακέτα εγκατάστασης συστή$\mu$ατος συχνά χρησι$\mu$οποιού$\mu$ενα. & αποσυ$\mu$πίεση App/Virtual drive \tabularnewline
HTML & Ιστοσελίδα & Αρχείο ιστού, ανοίγει $\mu$ε browser. Σε αποθηκευ$\mu$ένες ιστοσελίδες συχνά χρησι$\mu$οποιού$\mu$ενα. & Browser \tabularnewline
\hline
\end{tabularx}
\end{center}

\begin{tcbpractice}{Εξάσκηση: Οργανώστε τον φάκελο «Λήψεις»}
Ανοίξτε τον διαχειριστή αρχείων > $\mu$πείτε στον φάκελο «Λήψεις». Δείτε τι υπάρχει; Διαγράψτε τα άχρηστα, τα χρήσι$\mu$α ταξινο$\mu$ήστε (π.χ. δη$\mu$ιουργήστε φακέλους «Εργασία», «Memes»). Ελέγξτε τον αποθηκευτικό χώρο, δείτε ποιοι τύποι αρχείων καταλα$\mu$βάνουν τα περισσότερα.

\end{tcbpractice}

\booksection{VIII. Διαχείριση εικόνων και συστή$\mu$ατος αρχείων: τακτοποιήστε τα περιουσιακά στοιχεία}

Στο προηγού$\mu$ενο τ$\mu$ή$\mu$α $\mu$άθα$\mu$ε «να αναγνωρίζου$\mu$ε φακέλους», αλλά η αναγνώριση δεν αρκεί. Σε αυτή την ενότητα θα $\mu$ιλήσου$\mu$ε για\textbf{Πώς να δη$\mu$ιουργήσετε ενεργά ένα $\mu$ακροπρόθεσ$\mu$ο πλαίσιο διαχείρισης}, ώστε τα αρχεία σας να βρίσκονται πάντα, να $\mu$ην χάνονται και να $\mu$ην παραβιάζονται.

\booksubsection{Κανόνες ταξινό$\mu$ησης και ονο$\mu$ατοθεσίας αρχείων}

Η ρίζα της ακαταστασίας είναι η έλλειψη κανόνων. Η δη$\mu$ιουργία κανόνων είναι απλή, το κλειδί είναι\textbf{Επι$\mu$ονή}.

\begin{itemize}
  \item \textbf{Ονο$\mu$ατοθεσία «χρόνος+γεγονός»}: το όνο$\mu$α του φακέλου περιλα$\mu$βάνει χρονική σή$\mu$ανση και περιγραφή γεγονότος, διευκολύνοντας την ανάκληση και αναζήτηση. Για παράδειγ$\mu$α\texttt{202606\_Ταξιδιωτικά έγγραφα},\texttt{Απολογισ$\mu$ός τέλους έτους 2025},\texttt{Γενέθλια παιδιού 2024}. Έτσι $\mu$πορείτε να δείτε το χρονοδιάγρα$\mu$$\mu$α και να ξέρετε το περιεχό$\mu$ενο $\mu$ε $\mu$ια $\mu$ατιά.
  \item \textbf{Σύστη$\mu$α ταξινό$\mu$ησης τριών επιπέδων}:
  \begin{itemize}
    \item \textbf{Επίπεδο 1}: κατά τύπο αρχείου (έγγραφα / εικόνες / βίντεο / ήχος)
    \item \textbf{Επίπεδο 2}: ανά σενάριο χρήσης (εργασία / ζωή / $\mu$άθηση)
    \item \textbf{Επίπεδο 3}: κατά συγκεκρι$\mu$ένο γεγονός ή έργο (όπως "2026Q2\_Έργο A")
  \end{itemize}
  \item \textbf{Μηνιαίο σχέδιο καθαρισ$\mu$ού}: Ορίστε $\mu$ια σταθερή η$\mu$ερο$\mu$ηνία κάθε $\mu$ήνα (π.χ. η$\mu$έρα $\mu$ισθοδοσίας ή 1 η του $\mu$ήνα) για οργάνωση του φακέλου "Λήψεις":\textbf{Αρχειοθετήστε τα χρήσι$\mu$α, διαγράψτε τα άχρηστα, $\mu$εταφέρετε τα πολύτι$\mu$α}. 12 $\mu$ικρές οργανώσεις το χρόνο είναι πολύ πιο εύκολες από $\mu$ία $\mu$εγάλη ετήσια γενική καθαριότητα.
\end{itemize}

\begin{tcbpractice}{Λίστα εργασία αυτού του $\mu$ήνα}
Διαλέξτε ένα ελεύθερο Σαββατοκύριακο, ακολουθήστε τη σειρά:

\begin{enumerate}
  \item Δη$\mu$ιουργήστε στον διαχειριστή αρχείων\texttt{2026Q2\_Έργο A}φάκελος εργασίας αυτού του $\mu$ήνα·
  \item Δείτε όλα τα αρχεία αυτής της εβδο$\mu$άδας στον κατάλογο «Λήψεις», κρατήστε τα χρήσι$\mu$α, αρχειοθέτηση, διαγράψτε τα άχρηστα·
  \item Δη$\mu$ιουργήστε για τις φετινές οικογενειακές φωτογραφίες\texttt{Οικογένεια 2025},\texttt{Οικογένεια 2026}τέτοιους ετήσιους φακέλους.
\end{enumerate}

\end{tcbpractice}

\booksubsection{Στρατηγική διαχείρισης περιττών αρχείων}

Το κινητό$\mu$ε παρατετα$\mu$ένη χρήση «παχαίνει», κυρίως λόγω περιττών αρχείων.

\begin{itemize}
  \item \textbf{Αυτό$\mu$ατο καθάρισ$\mu$α 72 ωρών}: προσωρινά αρχεία που λα$\mu$βάνονται από κοινωνικές εφαρ$\mu$ογές όπως WeChat, QQ (αυτοκόλλητα, εικόνες συνο$\mu$ιλίας, $\mu$ικρά βίντεο), τα περισσότερα είναι άχρηστα $\mu$ετά από $\mu$ία προβολή. Στο [WeChat] > [Ρυθ$\mu$ίσεις] > [Γενικά] > [Χώρος αποθήκευσης] ενεργοποιήστε τον "Αυτό$\mu$ατο καθαρισ$\mu$ό", ρυθ$\mu$ίστε τη ληγ$\mu$ένη προσωρινή $\mu$νή$\mu$η στις 72 ώρες. Το κινητό γίνεται α$\mu$έσως πιο καθαρό, χωρίς να χάνονται ση$\mu$αντικές εγγραφές.
  \item \textbf{Στοχευ$\mu$ένο καθάρισ$\mu$α χώρου αποθήκευσης}: Android [Ρυθ$\mu$ίσεις] > [Αποθήκευση] > [Εκκαθάριση], iOS [Ρυθ$\mu$ίσεις] > [Γενικά] > [Αποθήκευση iPhone], το σύστη$\mu$α ε$\mu$φανίζει ταξινό$\mu$ηση «cache / προσωρινά αρχεία / προς απεγκατάσταση App / $\mu$εγάλα αρχεία». Προτεραιότητα: (1) cache, (2) παλιά πακέτα εγκατάστασης (από αναβαθ$\mu$ίσεις), (3) διπλότυπα αρχεία (ίδιες φωτογραφίες).
  \item \textbf{Κανόνας αυτό$\mu$ατης ταξινό$\mu$ησης}: οι περισσότερες εφαρ$\mu$ογές άλ$\mu$που$\mu$ υποστηρίζουν έξυπνη ταξινό$\mu$ηση: πρόσωπα, τοπία, έγγραφα, στιγ$\mu$ιότυπα. Μετά την ενεργοποίηση, οι νέες ο$\mu$αδικές φωτογραφίες κατατάσσονται αυτό$\mu$ατα σε "Πρόσωπα > Όνο$\mu$α ", τα σαρω$\mu$ένα συ$\mu$βόλαια σε "Έγγραφα".\textbf{Σε συνδυασ$\mu$ό $\mu$ε τη $\mu$ηνιαία οργάνωση, $\mu$ειώνει ση$\mu$αντικά την πιθανότητα «να $\mu$ην βρίσκετε τα πράγ$\mu$ατά σας»}.
\end{itemize}

\booksubsection{Πρότυπο επεξεργασίας ασφάλειας εικόνων}

Οι εικόνες είναι ο πιο “ευαίσθητος” τύπος δεδο$\mu$ένων στο τηλέφωνο — $\mu$ια τυχαία λήψη $\mu$πορεί να αποκαλύψει την τοποθεσία, τη συσκευή και τον χρόνο σας. Συνιστού$\mu$ε\textbf{Δη$\mu$ιουργία ροής εργασίας “διπλής προστασίας”}:

\begin{tcbwarn}{Σταθερή διαδικασία αποστολής φωτογραφίες εγγράφων/ευαίσθητων screenshots}
\begin{enumerate}
  \item \textbf{Μαρκάρισ$\mu$α ορατών πληροφοριών}: χρησι$\mu$οποιήστε ενσω$\mu$ατω$\mu$ένα εργαλεία ή τρίτες εφαρ$\mu$ογές (π.χ. Meitu Xiuxiu, Snapseed)\textbf{Καλύψτε τα $\mu$εσαία ψηφία του αριθ$\mu$ού ταυτότητας, πινακίδα αυτοκινήτου, αριθ$\mu$ό πόρτας, όνο$\mu$α, ποσό}και άλλα κρίσι$\mu$α πεδία·
  \item \textbf{Αφαίρεση $\mu$εταδεδο$\mu$ένων}: πριν την αποστολή στο άλ$\mu$που$\mu$ πατήστε "Κοινοποίηση " ή "Εξαγωγή", το σύστη$\mu$α θα ρωτήσει αν θέλετε να διατηρήσετε τις πληροφορίες τοποθεσίας, επιλέξτε\textbf{"Μη διατήρηση"}. Έτσι το $\mu$οντέλο της συσκευής, οι συντεταγ$\mu$ένες GPS και η ώρα λήψης δεν αποστέλλονται $\mu$αζί $\mu$ε την εικόνα.
  \item \textbf{Προσωρινό watermark}(προαιρετικό): σε ευαίσθητα screenshots προσθέστε watermark κει$\mu$ένου «$\mu$όνο για χρήση XX», αποφεύγετε κατάχρηση.
\end{enumerate}

\end{tcbwarn}

:Συνηθισ$\mu$ένες $\mu$έθοδοι καθαρισ$\mu$ού $\mu$εταδεδο$\mu$ένων (EXIF) (βασικές λειτουργίες, για βαθύτερη προστασία δείτε Κεφάλαιο 6 «Αποευαισθητοποίηση φωτογραφιών»):

\begin{itemize}
  \item \textbf{Ενσω$\mu$ατω$\mu$ένο σύστη$\mu$α}: iOS [Ρυθ$\mu$ίσεις] > [Απόρρητο και ασφάλεια] > [Κά$\mu$ερα], $\mu$πορείτε να απενεργοποιήσετε την "Ακριβή τοποθεσία "· κατά την κοινοποίηση επιλέξτε "Μόνο τρέχουσα φωτογραφία" για να\textbf{Αυτό$\mu$ατη αφαίρεση των περισσότερων EXIF}.
  \item \textbf{Android}: Οι περισσότερες εφαρ$\mu$ογές γκαλερί $\mu$άρκας ζητούν "Αφαίρεση πληροφοριών τοποθεσίας" κατά την κοινοποίηση — απλά επιλέξτε το.
  \item \textbf{Επαγγελ$\mu$ατικά εργαλεία}: Το ExifTool (υπολογιστής), Metapho (iOS), Photo Metadata Remover (Android) $\mu$πορούν να αφαιρέσουν πλήρως όλα τα $\mu$εταδεδο$\mu$ένα.
\end{itemize}

Για όσους χρειάζονται\textbf{Μακροπρόθεσ$\mu$η αποθήκευση}ευαίσθητες εικόνες (ταυτότητα, τραπεζική κάρτα, ληξιαρχικές πράξεις, διαβατήριο), συνιστού$\mu$ε να αποθηκεύονται σε\textbf{Κρυπτογραφη$\mu$ένο άλ$\mu$που$\mu$}:

\begin{itemize}
  \item iPhone: [Φωτογραφίες] > επιλέξτε φωτογραφία > Απόκρυψη (προβολή $\mu$όνο $\mu$ετά από Face ID επικύρωση)·
  \item Huawei/Xiaomi: [Άλ$\mu$που$\mu$] > long-press φωτογραφία > «Μετακίνηση στο χρη$\mu$ατοκιβώτιο»·
  \item Εφαρ$\mu$ογές κρυπτογράφησης τρίτων:\textbf{360 Security Album},\textbf{Tencent Album Manager (ιδιωτικό άλ$\mu$που$\mu$)},\textbf{KeepSafe}κ.λπ., $\mu$ε ξεχωριστή προστασία κωδικού — καλύτερη $\mu$υστικότητα.
\end{itemize}
